\documentclass{article}
\usepackage{amsmath}

\title{A Minimal Attention Model}
\author{Test Author}

\begin{document}
\maketitle

\begin{abstract}
We describe a small attention model used as an ingestion fixture.
\end{abstract}

\section{Introduction}
Recurrent models process sequences step by step, which limits parallelism.
We propose an attention-only architecture following \citep{bahdanau2014}.

% This comment must never appear in the parsed output.
The model is \emph{simple} and \textbf{fast} to train.

\section{Method}
\label{sec:method}

Attention is computed over queries, keys, and values:

\begin{equation}
\label{eq:attention}
\mathrm{Attention}(Q,K,V) = \mathrm{softmax}\!\left(\frac{QK^{\top}}{\sqrt{d_k}}\right)V
\end{equation}

where $d_k$ is the key dimension. See Equation~\ref{eq:attention} for details.

\begin{algorithm}
\caption{Training loop}
\label{alg:train}
\begin{algorithmic}
\STATE Initialize parameters
\STATE Sample a minibatch
\end{algorithmic}
\end{algorithm}

\section{Experiments}
We train with the Adam optimizer using a learning rate of 0.0001 and a dropout
rate of 0.1 on the base model. Training runs for 100000 steps.

\begin{figure}
\includegraphics{plot.png}
\caption{A figure whose body should not become prose.}
\end{figure}

\begin{thebibliography}{9}
\bibitem{bahdanau2014} Bahdanau et al. Neural machine translation. arXiv:1409.0473, 2014.
\end{thebibliography}

\end{document}
